\scanpage{118}
\mspage{58}
et il faut demander de plus qu'on ait commutativité dans le \deleted{\unsure{simple}} diagramme

(177)\quad
\diagram{\begin{tikzcd}[ampersand replacement=\&, column sep=large,row sep=large]
D_{d,d',d''}\arrow[r]\arrow[d,"\mathrm{tautol.}"'] \& D_{d',d''}\arrow[d,"\text{``but''}"]\\
D_{d,d''}\arrow[r,"\text{``but''}"'] \& D_{d''}
\end{tikzcd}}.

\uline{Remarque} L'axiomatique précédente a un sens dans chaque site (\deleted{\unsure{ce qui}} le fait de travailler dans un site donne un sens à la notion de revêtement principal\ldots), où on suppose que les fonctions loc constantes définissent une décomposition en sommes disjointes d'un objet. Il est inutile en fait de se \uline{donner} la fonction codimension, elle est contenue dans la condition c), qui implique a) et b).

Pour faire le lien avec l'intuition ``croisements normaux'', on va faire de plus une hypothèse ad hoc sur la nature des morphismes
\[
D_{d,d'}\longrightarrow D_{d'}
\]
(NB $D_{d,d'}\longrightarrow D_d$ est un rev. fini étale\ldots), p.ex.
\margin{117}
